\chapter{Data quality: gaps and completeness}\label{ch:quality}

A seismograph is only useful while it is recording. The application measures and displays three aspects of
\emph{data availability}: what is on the laptop and what is still on the Shake (the coverage bar), what is missing
inside the selected window (Data health), and whether the Shake's clock can be trusted (Chapter~\ref{ch:data}).

\section{Samples, spans and gaps}
A trace is a set of samples at times $t_i=t_0+i/\fs$. Records and traces end where a gap begins, so the whole day is a
list of \emph{spans} $[s_j,e_j)$, where $e_j$ is the time just after the last sample of the span, $e_j=s_j+N_j/\fs$. Two
spans are \emph{contiguous} when the hole between them is at most half a sample, which absorbs timing rounding;
contiguous spans are merged. What counts as a gap depends on who is asking, so the application uses three tolerances.

\begin{table}[tbp]
\centering\small
\caption{Gap tolerances used by different parts of the application.}\label{tab:tol}
\begin{tabular}{@{}L{0.34\linewidth}cL{0.46\linewidth}@{}}
\toprule
Use & Tolerance & Rule\\ \midrule
Coverage bar (\file{core/coverage.py}) & $0.5/\fs$ & spans closer than half a sample are merged into one\\
Data health (\file{core/health.py}) & $0.5/\fs$ & same rule, applied per channel\\
Helicorder rows (\file{core/helicorder.py}) & $1.5/\fs$ & a break is drawn only when the hole exceeds one and a
  half sample intervals\\
\bottomrule
\end{tabular}
\end{table}

\section{Completeness}
For a selection $[a,b)$ with $\Tsel=b-a$ the application computes, for every channel, the set of times with data (the union
of its clipped spans) and from it the set of times \emph{missing} on that channel. The \emph{gaps} of the selection are the
union, over channels, of the missing sets: a moment counts as a gap if \emph{any} selected channel has no data. The
\emph{completeness} is
\begin{equation}\label{eq:complete}
  C = 1-\frac{\sum_g (e_g - s_g)}{\Tsel}, \qquad 0\le C\le1,
\end{equation}
where the sum runs over the gaps $g$ of the union. If no channel has any data in the selection there is one gap equal to
the whole selection and $C=0$. \Cref{fig:gaps} illustrates the construction.

\begin{figure}[tbp]
\centering
\begin{tikzpicture}[x=0.078\linewidth,y=1cm]
  \foreach \y/\lab/\spans in {2.4/EHZ/{0/3.2, 4.0/10},1.6/EHN/{0/10},0.8/EHE/{0/6.1, 6.7/10}} {
    \node[left,font=\small\sffamily] at (-0.15,\y) {\lab};
    \foreach \a/\b in \spans { \fill[rs3dblue!70] (\a,\y-0.22) rectangle (\b,\y+0.22); }
  }
  \node[left,font=\small\sffamily] at (-0.15,0) {union};
  \fill[rs3dred!70] (3.2,-0.22) rectangle (4.0,0.22);
  \fill[rs3dred!70] (6.1,-0.22) rectangle (6.7,0.22);
  \draw[rs3dgrey] (0,-0.6) -- (10,-0.6);
  \foreach \x/\t in {0/$a$,10/$b$} {\draw[rs3dgrey] (\x,-0.5) -- (\x,-0.7); \node[below,font=\small] at (\x,-0.7) {\t};}
  \node[font=\small,rs3dred!80!black,above] at (3.6,0.25) {$g_1$};
  \node[font=\small,rs3dred!80!black,above] at (6.4,0.25) {$g_2$};
  \draw[decorate,decoration={brace,amplitude=3pt,mirror},rs3dgrey] (0,-1.15) -- (10,-1.15)
    node[midway,below=3pt,font=\small] {$\Tsel$};
\end{tikzpicture}
\caption{Completeness of a selection. Each channel has its own holes; the gaps of the selection ($g_1$, $g_2$) are the
union of all of them, and $C=1-(|g_1|+|g_2|)/\Tsel$. In this sketch $C = 1-(0.8+0.6)/10 = 86\,\%$.}\label{fig:gaps}
\end{figure}

\begin{rs3dapp}
\code{health(stream, start, end)} in \file{core/health.py} returns the completeness and the list of gaps, each as a
pair of times; the Data page prints it as a percentage and lists up to 50 gaps with their durations in the chosen
time zone. The leading and trailing holes of a selection (data starting after $a$ or ending before $b$) are gaps like
any other.
\end{rs3dapp}

\section{Coverage: laptop, Shake, or nothing}
The coverage bar above the Data page's preview classifies every instant of a view into three states:
\emph{local} (the samples are on the laptop), \emph{remote} (the day exists on the Shake but is not downloaded) and
\emph{gap} (no data anywhere). The rule, evaluated per segment between cut points at midnight and at span edges, is:
\begin{enumerate}
\item if the midpoint of the segment lies inside a local span, the state is local;
\item otherwise, if that day exists on the Shake \emph{and} the laptop does not have the day file, or the day is the newest day
  on the Shake (still being written), the state is remote;
\item otherwise it is a gap.
\end{enumerate}
The ``newest day'' exception is why today's file is always offered for download again: the Shake appends to it
continuously, so a copy made earlier is incomplete.

\section{Why a gap is never filled}
Interpolating across a gap would be easy and would make plots prettier, but it would manufacture data. RS3D keeps gaps as
gaps in every stage: separate traces on loading, separate filter runs, \code{NaN} separators in the plotted lines, blank
spectrogram columns, and an explicit completeness figure. When a gap is shorter than the window of an operation, that
operation excludes the window rather than guessing its content.
